% Shaurya's section, transcribed verbatim from his compiled report (report-1.pdf).
\section{Introduction}
The physical layout of a table decides which bytes a query moves through the memory hierarchy.
NSM stores each tuple contiguously, so every query reads whole tuples. DSM stores each attribute
contiguously, so a query reads only the attributes it needs but has to stitch them back into
tuples. PAX~\cite{pax} keeps NSM's page but stores each attribute of the page in its own
\emph{minipage}, aiming for DSM's cache behaviour with NSM's locality.

\textbf{Research question.} \emph{How do the layout and the materialization strategy affect the
latency and cache behaviour of analytical scans across projectivity, selectivity, working-set
size and PAX page size, and how should the code change to run faster on one specific machine?}
We study four crossovers: how many attributes before NSM beats DSM, how many rows before early
beats late materialization~\cite{abadi}, whether the answers change when the table fits in
cache, and how large a PAX page must be to scan like DSM. For the machine-specific half we
start, as the project suggests, with materialization (which lines are touched) and compression
(how many values fit in a line), then extend to vector size, memory placement, output writes
and page size.

Sections 2 and 3 describe the engine and the measurements. Section 4 explains the layout
results by counting cache lines, Sections 5 and 6 apply the count to materialization and
encodings, and Section 7 turns the machine's parameters into tuning rules.
